% Pandoc template for the fed_fraud documents.
% Based on proposal_gen/templates/detailed_proposal_template/detailed_1_classic.tex
% (Classic report): centered title page, table of contents, ruled section
% headings and the Geodesic house footer. Fonts are Latin Modern (lmodern).
% Two layouts: the full report (title page and contents) and, with
% -V brief, a compact one-page-brief header in the same style.
\documentclass[letterpaper,11pt]{article}
\usepackage[T1]{fontenc}
\usepackage[utf8]{inputenc}
\usepackage{lmodern}
\usepackage[letterpaper,top=0.9in,bottom=1.2in,left=1in,right=1in,footskip=0.6in]{geometry}
\usepackage[table]{xcolor}
\usepackage{geodesicbrand}
\usepackage{amsmath,amssymb}
\usepackage{enumitem}
\usepackage{array}
\usepackage{calc}
\usepackage{longtable}
\usepackage{makecell}
\usepackage[most]{tcolorbox}
\usepackage{titlesec}
\usepackage{needspace}
\usepackage{etoolbox}
\usepackage{hyperref}
\hypersetup{colorlinks=true,urlcolor=coolblack,linkcolor=coolblack,filecolor=geoblue,
  pdftitle={Office of Personnel Management},pdfauthor={\geofirm}}
\urlstyle{same}

\setlist{itemsep=2pt,leftmargin=1.4em}
\setlength{\parindent}{0pt}
\setlength{\parskip}{6pt}
\providecommand{\tightlist}{\setlength{\itemsep}{0pt}\setlength{\parskip}{0pt}}
\setcounter{secnumdepth}{0}

% Section style: house rule under each heading
\titleformat{\section}{\Large\bfseries\color{coolblack}}{\thesection}{0.6em}{}[{\color{coolblack}\titlerule[0.5pt]}]
\titleformat{\subsection}{\large\bfseries\color{geomid}}{\thesubsection}{0.6em}{}
\titleformat{\subsubsection}{\normalsize\bfseries\color{geomid}}{\thesubsubsection}{0.6em}{}
\titlespacing*{\section}{0pt}{10pt}{6pt}

% Tables are written by latex/filters.lua (navy header, banded rows)
\setlength{\LTpre}{6pt}\setlength{\LTpost}{6pt}

\begin{document}

% BRIEF HEADER
\thispagestyle{geodesic}
\begin{minipage}[c]{0.13\textwidth}\includegraphics[width=\linewidth]{icon.png}\end{minipage}\hfill
\begin{minipage}[c]{0.84\textwidth}
  {\textsc{\large \geofirm}}\par\vspace{2pt}
  {\small\color{geogray}Federal fraud expansion\enspace\textperiodcentered\enspace When: \compiletime}
\end{minipage}\par\vspace{6pt}
{\color{coolblack}\rule{\linewidth}{0.3mm}}\par\vspace{2pt}
{\LARGE\bfseries\color{coolblack} Office of Personnel Management}\par
\vspace{2pt}{\large\color{geomid} One-page brief}\par
\vspace{-2pt}{\color{coolblack}\rule{\linewidth}{0.3mm}}\par\vspace{4pt}
\pagestyle{geodesic}

\textbf{Priority:} 11 of 33 (score 3.3). \textbf{FY2026 outlays:} \$201B
(USAspending). \textbf{As of:} 2026-10-27. Public data only.

\hypertarget{bottom-line}{%
\section{Bottom line}\label{bottom-line}}

OPM runs two big payment streams: federal retirement (\$109B in FY2025)
and the Federal Employees Health Benefits program (about \$70B, 8.2
million members). Retirement error rates are low, but GAO found in April
2026 that FEHB paid deceased, incarcerated and excluded providers, and
that OPM's own debarment list is out of step with the HHS and SAM
exclusion lists. That cross-list gap is something we can measure with
public data alone.

\hypertarget{exposure}{%
\section{Exposure}\label{exposure}}

\begin{itemize}
\tightlist
\item
  Federal Retirement Services: \$109.4B outlays, about \$364M improper
  (0.33\%) in FY2025.
\item
  FEHB: about 400 deceased providers paid for services at least 50 days
  after death; over 2,000 claims from providers excluded elsewhere; over
  2,500 claims from providers OPM had already debarred. (GAO, CY2022
  claims)
\item
  Over 85\% of OPM OIG's debarment records lack an NPI, and at least
  1,200 LEIE entries that meet OPM's debarment criteria are missing from
  its list.
\end{itemize}

\hypertarget{public-data-we-can-use}{%
\section{Public data we can use}\label{public-data-we-can-use}}

\begin{itemize}
\tightlist
\item
  \textbf{HHS-OIG LEIE, SAM.gov exclusions, and OPM OIG's debarment
  list.}
\item
  \textbf{NPPES} NPI registry (addresses, taxonomies, deactivation
  dates).
\item
  \textbf{CMS provider data} to profile FEHB-overlapping providers.
\item
  \textbf{OPM OIG investigative summaries} as labeled cases.
\end{itemize}

\hypertarget{detection-plays}{%
\section{Detection plays}\label{detection-plays}}

\begin{enumerate}
\def\labelenumi{\arabic{enumi}.}
\tightlist
\item
  \textbf{Exclusion list reconciliation:} match the three exclusion
  lists on name, NPI, address and license to find providers barred in
  one program and not another.
\item
  \textbf{Dead and deactivated providers:} NPIs deactivated or tied to
  deceased individuals that still look active in carrier networks.
\item
  \textbf{Cross-program spillover:} providers flagged in Medicare work
  (see \url{hhs.pdf}) who also serve FEHB members.
\end{enumerate}

\hypertarget{first-prototype}{%
\section{First prototype}\label{first-prototype}}

An exclusion-list reconciliation report: every provider excluded by HHS
or SAM but missing from OPM's list, with matched identifiers. This
reuses the HHS provider pipeline almost unchanged.

\hypertarget{limits}{%
\section{Limits}\label{limits}}

FEHB claims are held by carriers and are not public, so we cannot see
what was paid to whom. Retirement payee data is not public.

\hypertarget{who-to-take-it-to}{%
\section{Who to take it to}\label{who-to-take-it-to}}

OPM Healthcare and Insurance, OPM OIG, and FEHB carriers' special
investigation units.

\hypertarget{sources}{%
\section{Sources}\label{sources}}

\begin{itemize}
\tightlist
\item
  \href{https://files.gao.gov/reports/GAO-26-108139/index.html}{GAO-26-108139:
  FEHB provider eligibility risks (April 2026)}
\item
  \href{https://paymentaccuracy.gov/program/opm-federal-retirement-services}{PaymentAccuracy.gov:
  OPM Federal Retirement Services}
\item
  \href{https://www.fedweek.com/issue-briefs/bogus-claims-cashing-checks-of-deceased-ig-cites-cases-of-fehbp-retirement-fraud/}{FedWeek:
  OPM OIG retirement and FEHB fraud cases}
\item
  \href{https://www.usaspending.gov/agency}{USAspending agency profiles
  (FY2026 outlays)}
\end{itemize}

\end{document}
